%AAAAAAAAAAAAAAAAAAAAAAAAAAAAAAAAAAAAAAAAAAAAAAAAAAAAAAAAAAAAAAAAAAAAAAAAAAAAAAAAAAAAAAAAAAAAAAAAAAAAAAAAAAAAAAA
%AAAAAAAAAAAAAAAAAAAAAAAAAAAAAAAAAAAAAAAAAAAAAAAAAAAAAAAAAAAAAAAAAAAAAAAAAAAAAAAAAAAAAAAAAAAAAAAAAAAAAAAAAAAAAAA

%Ce document est r�alis� par
%Mlle S. BOUCHELAGHEM Doctorante/Enseignante
%D�partement d'Informatique, Universit� de B�jaia

%Pour les �tudiants de Licence 3 en Informatique
%Dans le cadre du module *** R�daction Scientifique ***

% ***Avril 2018***

%AAAAAAAAAAAAAAAAAAAAAAAAAAAAAAAAAAAAAAAAAAAAAAAAAAAAAAAAAAAAAAAAAAAAAAAAAAAAAAAAAAAAAAAAAAAAAAAAAAAAAAAAAAAAAAA
%AAAAAAAAAAAAAAAAAAAAAAAAAAAAAAAAAAAAAAAAAAAAAAAAAAAAAAAAAAAAAAAAAAAAAAAAAAAAAAAAAAAAAAAAAAAAAAAAAAAAAAAAAAAAAAA

\documentclass[12pt]{report}
\usepackage[latin1]{inputenc}
\usepackage[T1]{fontenc}
\usepackage[francais]{babel}
\usepackage{graphicx}
\usepackage{layout}
\usepackage[top=2cm, bottom=2cm, left=2cm, right=2cm]{geometry}
\usepackage{url}
\usepackage{fancyhdr}
\usepackage{color}
\usepackage{colortbl}
\usepackage{multirow}
\usepackage{nomencl}
\usepackage{setspace}


\makenomenclature

\definecolor{lightgray}{gray}{0.85}
\pagestyle{fancy}

\fancyhead[L]{}

\begin{document} 

%%%% Page de garde %%%%
%%%%%%%%%%%%%%%%%%%%%%%
%%%% Page de garde %%%%
%%%%%%%%%%%%%%%%%%%%%%%

\begin{titlepage}

\begin{center}

  \large{R�publique Alg�rienne D�mocratique et Populaire}\\
  \large{Minist�re de l'Enseignement Sup�rieur et de la Recherche Scientifique}\\
  \large{Universit� A. Mira de B�ja�a}\\
  \large{Facult� des Sciences Exactes}\\
  \large{D�partement d'Informatique}\\


	\begin{figure}[h!]
		\centering
		\includegraphics[width=6cm]{images/logo.jpg}
	\end{figure}
	
    \paragraph{}
     
		{\LARGE{\textbf M�moire de Fin de Cycle}} \\[2ex]
    \textbf En vue de l'obtention du dipl�me de Licence en Informatique G�n�rale \\[1ex]
    \paragraph{}
		
    \textbf{Th�me}\\[1ex]
		\rule{18cm}{1pt}
        {\Huge{\textbf{{\\
    Intitul� de votre m�moire \\}}}}
		\rule{18cm}{1pt}\\
	  \vspace{0.2cm}

\end{center}

\begin{center}

\large R�alis� par

\paragraph{}

\begin{tabular}{l l l l}	

		M. NOM Pr�nom 					& & & 		M. NOM Pr�nom \\
								
		M$^{lle}$ NOM Pr�nom 		& & & 		M. NOM Pr�nom \\		
		
		M$^{lle}$ NOM Pr�nom 		& & & 		M$^{lle}$ NOM Pr�nom \\		
		
\end{tabular}
		
\paragraph{}

\large Devant le jury compos� de

\end{center}
  
\paragraph{}
  
\begin{table}[h!]
\begin{center}
        \begin{tabular}{l l l l}
				\textbf{Pr�sident :} & M. NOM Pr�nom & Grade de l'enseignant & Universit� de B�ja�a \\
				
				\textbf{Examinateur :} & M. NOM Pr�nom & Grade de l'enseignant & Universit� de B�ja�a \\
				
				\textbf{Examinatrice :} & M$^{me}$ NOM Pr�nom & Grade de l'enseignante & Universit� de B�ja�a \\
				
				\textbf{Encadrant :} & M. NOM Pr�nom & Grade de l'enseignant & Universit� de B�ja�a \\
				
        \end{tabular}    
\end{center}

\end{table}

\paragraph{}

\begin{center} Promotion 2017 - 2018 \end{center}


\end{titlepage}


%%%% Remerciements %%%%
\begin{titlepage}
\newpage
\pagestyle{fancy}      
\lhead{}  
\chead{}     
\rhead{}     
    
\renewcommand{\headrulewidth}{0.5pt}

\chapter*{\hrulefill ~~\textbf{Remerciements}~ \hrulefill}
\paragraph{}
\paragraph{}
\large

	  \textbf{T}exte ...
		
		\paragraph{}
			
		\textbf{T}exte ...
		
		\paragraph{}
		
		
		\textbf{T}exte ...
		
		\paragraph{}
		
		\textbf{T}exte ...
		
		\paragraph{}
		
		\textbf{T}exte ...

\thispagestyle{empty}		

\end{titlepage}	

%%%% D�dicaces %%%%
\input{D�dicaces}

%D�but de la num�rotation romaine pour la table des mati�res
\newpage
\pagenumbering{roman}

%Construire la table des mati�res
\addcontentsline{toc}{part}{Table des mati�res} 
\tableofcontents
\fancyhead[R]{}
\renewcommand{\headrulewidth}{0pt}

%Construire la table des figures
\listoffigures
\addcontentsline{toc}{part}{Table des figures}

%Construire la liste des tableaux
\listoftables
\addcontentsline{toc}{part}{Liste des tableaux}

%Construire la liste des abr�viations
\nomenclature{\textbf{UML}}{Unified Modeling Language}
\nomenclature{\textbf{UP}}{Unified Process}
\nomenclature{\textbf{SE}}{Syst�me d'Exploitation}


\def\nomname{Liste des abr�viations}
\printnomenclature[2in] 
\addcontentsline{toc}{part}{Liste des abr�viations}
\fancyhead[R]{\textit{Liste des abr�viations}}
\renewcommand{\headrulewidth}{1pt}

%D�but de la num�rotation arabe pour le reste du m�moire
\newpage
\pagenumbering{arabic}

%%%% Mise en forme des titres des chapitres %%%%
\makeatletter
\def\@makechapterhead#1{%
  \vspace*{10\p@}%
  {\parindent \z@ \centering \normalfont
    \ifnum \c@secnumdepth >\m@ne
        \huge\bfseries \@chapapp\space \thechapter
        \par\nobreak
        \vskip 20\p@
    \fi
    \interlinepenalty\@M
    \Huge \bfseries #1\par\nobreak
    \vskip 40\p@
  }}
 
\def\@makeschapterhead#1{%
  \vspace*{10\p@}%
  {\parindent \z@ \centering
    \normalfont
    \interlinepenalty\@M
    \Huge \bfseries  #1\par\nobreak
    \vskip 40\p@
  }}
\makeatother 

%%%% Profondeur des sous-titres jusqu'� 4 chiffres %%%%
\setcounter{secnumdepth}{3}

%%%% Introduction G�n�rale %%%%
\input{Introduction_G�n�rale}

%%%% Chapitre 1 %%%%
\chapter{Titre du premier chapitre}
\fancyhead[R]{\textit{Titre du premier chapitre}}
\renewcommand{\headrulewidth}{1pt}


\section{Introduction}

Texte ...


\section{Titre de la section}

Texte ...


\subsection{Titre de la premi�re sous-section}

Texte ... \cite{1}.


\subsection{Titre de la deuxi�me sous-section}

Texte ... \cite{2}.


\subsubsection{Titre de la premi�re sous-sous-section}

	Texte ... \\

\begin{itemize}

\item[\textbullet] Premier point;

\item[\textbullet] Deuxi�me point;

\item[\textbullet] Dernier point.\\

\end{itemize}


\subsubsection{Titre de la deuxi�me sous-sous-section}

Texte ...

    
\section{Conclusion}

Texte ...

%%%% Chapitre 2 %%%%
\chapter{Titre du deuxi�me chapitre}
\fancyhead[R]{\textit{Titre du deuxi�me chapitre}}
\renewcommand{\headrulewidth}{1pt}


\section{Introduction}

Texte ...


\section{Titre de la section}

Texte ...


\subsection{Titre de la premi�re sous-section}

La figure \ref{fig1} ...

\begin{figure}[h!]
	\centering
		\includegraphics[width=8cm]{images/MySQL.png}
	\caption{Titre de la figure.}
	\label{fig1}
\end{figure}


\subsection{Titre de la deuxi�me sous-section}

Texte ...


%%%% Ajout d'un tableau %%%%
\begin{table}[!h]
		\small
		\centering
		\footnotesize{
			\begin{tabular}{|p{4cm}|p{4cm}|p{4cm}|} %%% La taille des trois colonnes est �gale � 4cm %%%
				\hline
				\rowcolor{lightgray} \textbf{Colonne 1} & \textbf{Colonne 2} & \textbf{Colonne 3} \\
				\hline
				Ligne 1 Colonne 1 & Ligne 1 Colonne 2 & Ligne 1 Colonne 3 \\
				\hline
				Ligne 2 Colonne 1	& etc. 							& etc. \\
				\hline
			\end{tabular}}
		\caption{Titre du tableau.}	
\end{table}
	

\section{Conclusion}

Texte ...

%%%% Conclusion G�n�rale %%%%
\input{Conclusion_G�n�rale}

%%%% Bibliographie %%%%
\begin{thebibliography}{2}
\addcontentsline{toc}{part}{Bibliographie}

		\bibitem[1]{1} P. \textsc{Roques}, \emph{UML 2 Mod�liser une Application Web}, 4$^e$ �dition, Eyrolles, Paris, 2008.
		
		\bibitem[2]{2} Visual Paradigm, \emph{Leading UML, BPMN, EA, Agile and Project Management Software}, https://www.visual-paradigm.com/, Version 15.0, Consult� le 17/03/2018.
		
		\bibitem[3]{3} ...
		
\end{thebibliography}

%%%% R�sum�/Abstract %%%%
\input{R�sum�}

\end{document}